(1) 没有电场，粒子总动能 \(mv^2/2\) 守恒；磁矩不变给 \(v_\perp^2(s)=v^2\sin^2\alpha_0 B(s)/B_0\)。因而 \(v_\parallel^2(s)=v^2[1-\sin^2\alpha_0(1+s^2/L^2)]=v^2\cos^2\alpha_0-(v^2\sin^2\alpha_0/L^2)s^2\)。反射点满足 \(v_\parallel=0\)，即 \(s_m=L\cot\alpha_0\)，须注意镜点不意味着总速度为零。

(2) 以平行运动 \(ds/dt=v_\parallel\) 的能量式微分，得到 \(\ddot s=-(v^2\sin^2\alpha_0/L^2)s\)。于是导引中心沿 \(s\) 作频率 \(\omega_b=v\sin\alpha_0/L\) 的简谐反弹运动，若从 \(s=0\) 向正方向出发，则 \(s(t)=s_m\sin(\omega_b t)\)。完整两镜点往返周期 \(T_b=2\pi/\omega_b=2\pi L/(v\sin\alpha_0)\)，此处周期来源于磁场的特定二次剖面，不能套用任意磁镜。

(3) 镜点磁场 \(B_m=B_0[1+s_m^2/L^2]=B_0(1+\cot^2\alpha_0)=B_0/\sin^2\alpha_0\)。回旋角频率大小 \(\omega_c(s)=|q|B(s)/m\)，在镜点放大为 \(\omega_{c,m}=\omega_{c,0}/\sin^2\alpha_0\)。磁场不对带电粒子做功，动能在平行和垂直分量间转移；在镜点全部转为垂直回旋运动。

(4) 当 \(\alpha_0\to0^+\)，\(s_m\to\infty,T_b\to\infty\)；这一极限不仅是角度上的奇异，也会使轨道探索磁场模型很远处，从而破坏给定局部二次剖面的可行性。第一绝热不变量需要 \(\rho_L|\nabla B|/B\ll1\)，且 \(\omega_b\ll\omega_c\) 才能清楚地区分回旋与反弹时标。若碰撞或强场梯度破坏该条件，以上镜点关系不再保证精确。
